\documentclass[12pt, a4paper]{article}

% ---- Packages ----
\usepackage[utf8]{inputenc}
\usepackage[T1]{fontenc}
\usepackage{amsmath, amssymb, amsthm}
\usepackage{mathtools}
\usepackage{braket}       % Dirac notation
\usepackage{hyperref}
\usepackage[margin=1in]{geometry}
\usepackage{enumitem}
\usepackage{tikz-cd}      % Commutative diagrams (category theory)

% ---- Theorem environments ----
\newtheorem{axiom}{Axiom}
\newtheorem{definition}{Definition}
\newtheorem{proposition}{Proposition}
\newtheorem{conjecture}{Conjecture}
\newtheorem{remark}{Remark}

% ---- Title ----
\title{\textbf{The Relational Witnessing Framework:} \\ 
The Emergence of Time and Experience from Entanglement Entropy}
\author{Trevor Scott}
\date{March 2026}

\begin{document}
\maketitle

% ====================================================================
\begin{abstract}
This paper outlines a formal mathematical framework locating the emergence of time and subjective experience within a globally static, timeless quantum state. Synthesizing the Wheeler--DeWitt equation, Landauer's Principle, the Page--Wootters mechanism, and the Margolus--Levitin quantum speed limit, we propose that witnessing, entropy, and time are co-generated by a single relational structure---the partial trace across a bipartition of a timeless whole---and constitute three projections of the same underlying epistemic cut. Crucially, this framework does not propose new physics; rather, it defines the structural scaffold that any future theory of experience must map from. We formalize the epistemic nature of the partial trace, derive a relational time bound under an explicit thermodynamic optimality assumption, construct a flag complex encoding the entanglement topology of a witnessing subsystem, and define the categorical and topological constraints that any solution to the hard problem of consciousness must mathematically satisfy.

\medskip
\noindent \textbf{Methodological note.} Axioms I--IV constitute the deductive core of the framework, relying on established physics and information theory. Axiom V and the Topological Hypothesis (Section~\ref{sec:topo}) are conjectures about the form that any bridge between physics and phenomenology must take. We are explicit about where derivation ends and conjecture begins.
\end{abstract}

% ====================================================================
\section{Axiom I: The Timeless Whole}
\label{sec:axiom1}

We begin outside of time. The universe, in its entirety, is a closed system in a pure, static quantum state. There is no external clock, and therefore no objective temporal evolution.

Let the total Hilbert space of the universe be $\mathcal{H}_{\mathrm{total}}$. The universe is described by a single, pure state vector $\ket{\Psi}$. The Hamiltonian $\hat{H}$ of the entire universe yields zero energy, meaning the state is strictly invariant under global time translation \cite{DeWitt1967}:
\begin{axiom}[Global Stasis]
\begin{equation}
\hat{H}\ket{\Psi} = 0
\label{eq:WDW}
\end{equation}
\end{axiom}

\noindent Because $\ket{\Psi}$ is a globally pure state, its density matrix $\rho_{\mathrm{total}} = \ket{\Psi}\!\bra{\Psi}$ has von Neumann entropy of exactly zero:
\begin{equation}
S(\rho_{\mathrm{total}}) = -\operatorname{Tr}(\rho_{\mathrm{total}} \ln \rho_{\mathrm{total}}) = 0.
\label{eq:Stotal}
\end{equation}

There is no arrow of time at the global level.

% ====================================================================
\section{Axiom II: The Hilbert Space Bipartition}
\label{sec:axiom2}

To yield the emergence of relational dynamics, the timeless whole must be partitioned. We mathematically divide the whole into a \emph{Witness} ($W$) and the \emph{Rest of the Universe} ($R$) by factoring the Hilbert space via a tensor product decomposition:

\begin{axiom}[Bipartition]
\begin{equation}
\mathcal{H}_{\mathrm{total}} = \mathcal{H}_W \otimes \mathcal{H}_R
\label{eq:bipartition}
\end{equation}
\end{axiom}

As noted by Zanardi, Lidar, and Lloyd \cite{Zanardi2004}, the tensor product structure of a Hilbert space is not uniquely given \emph{a priori}; it is observable-induced. We adopt the following criterion for the preferred factorization:

\begin{definition}[Preferred Factorization]
The dynamically preferred bipartition $\mathcal{H}_W \otimes \mathcal{H}_R$ is the tensor product structure with respect to which a specified algebra of observables $\mathcal{A}_W$ acts locally on $\mathcal{H}_W$.
\end{definition}

This is a \emph{static}, algebraic criterion: it selects the factorization by reference to the structure of the observable algebra, not by reference to any dynamical process. We note that Zurek's Quantum Darwinism \cite{Zurek2009} provides a dynamical confirmation of this selection in the emergent-time regime---the environment selectively proliferates information about pointer states, stabilizing precisely the factorization that makes $\mathcal{A}_W$ local. However, Quantum Darwinism is not invoked foundationally here, since it presupposes time, which is not defined until Axiom~IV. Instead, it serves as a \emph{consistency check}: the factorization selected algebraically must coincide with the one stabilized dynamically once relational time has emerged.

% ====================================================================
\section{Axiom III: The Epistemic Bipartition}
\label{sec:axiom3}

\begin{axiom}[Formal Witness]
A subsystem $W$ acts as a \emph{formal Witness} if and only if it possesses genuine quantum entanglement with $R$, characterized by:
\begin{enumerate}[label=(\roman*)]
    \item a non-trivial Schmidt rank in the decomposition of $\ket{\Psi}$ across the bipartition, and
    \item non-zero logarithmic negativity $E_{\mathcal{N}}(W{:}R) > 0$.
\end{enumerate}
\end{axiom}

The requirement of non-zero logarithmic negativity distinguishes genuine quantum entanglement from classical thermodynamic mixedness. A subsystem whose reduced state is mixed solely due to classical ignorance (a separable state) does not qualify as a formal Witness.

\begin{remark}
For the bipartition of a pure global state $\ket{\Psi}$, conditions (i) and (ii) are equivalent: non-trivial Schmidt rank implies non-zero logarithmic negativity, and vice versa. This follows because Schmidt rank greater than one for a pure state implies the reduced state $\rho_W$ is mixed, and for reductions of pure bipartite states, non-trivial mixedness implies positive partial-transpose negativity. Both conditions are stated explicitly because they diverge for \emph{mixed} bipartite states, which arise in the internal sub-factorizations of $W$ (Section~\ref{sec:topo}), where reduced states of subsystem pairs are generically mixed and can be mixed without being entangled.
\end{remark}

\begin{remark}
Whether a formal Witness possesses \emph{phenomenal experience} is a separate question entirely, addressed in Axiom~V.
\end{remark}

``Witnessing'' is not an ontological event that physically destroys superpositions; it is the epistemological adoption of a local perspective on a globally pure state. The Witness cannot access the entire universe $\mathcal{H}_{\mathrm{total}}$; it is restricted to its own local subsystem $\mathcal{H}_W$. To find the state of the Witness, we apply the partial trace, tracing out the rest of the universe:
\begin{equation}
\rho_W = \operatorname{Tr}_R\!\bigl(\ket{\Psi}\!\bra{\Psi}\bigr) = \sum_i p_i \ket{w_i}\!\bra{w_i}.
\label{eq:rhoW}
\end{equation}

Crucially, the partial trace does not describe what physically happens to the global wave function; it describes what the subsystem \emph{can know}. Due to this epistemic restriction, the von Neumann entropy of the Witness is strictly greater than zero:
\begin{equation}
S(\rho_W) = -\sum_i p_i \ln p_i > 0.
\label{eq:SW}
\end{equation}

Entropy is the intrinsic mathematical consequence of adopting a local perspective within an entangled system.

% ====================================================================
\section{Axiom IV: Relational Time}
\label{sec:axiom4}

Time is what entanglement entropy accumulation looks like from inside a bipartition. If the Witness acts as a ``clock,'' we define the state of the Rest of the universe conditional on the Witness reading a specific state $\ket{w(\tau)}$, following the Page--Wootters mechanism \cite{Page1983}\footnote{We use the standard shorthand $\bra{w(\tau)}\Psi\rangle$ for the rigorously defined conditional state $(\bra{w(\tau)}_W \otimes \mathbb{1}_R)\ket{\Psi}$. For the full perspective-neutral formalism, see H\"{o}hn, Smith, and Lock \cite{Hohn2021}.}:
\begin{equation}
\ket{\psi_R(\tau)} = \bigl(\bra{w(\tau)}_W \otimes \mathbb{1}_R\bigr)\ket{\Psi}.
\label{eq:PW}
\end{equation}

The parameter $\tau$ is the emergent, subjective time of the Witness. To anchor $\tau$ physically, we compose two established results under a limiting assumption.

\medskip
\noindent\textbf{Landauer's Principle} \cite{Landauer1961} establishes the minimum energy cost of erasing $\Delta S$ bits of information at temperature $T$:
\begin{equation}
\Delta E \geq k_B T \ln 2 \;\cdot\; \Delta S(\rho_W).
\label{eq:Landauer}
\end{equation}

\noindent\textbf{The Margolus--Levitin quantum speed limit} \cite{Margolus1998} bounds the minimum time for a quantum system to evolve to a distinguishable (orthogonal) state:
\begin{equation}
\Delta \tau \geq \frac{\pi \hbar}{2 \Delta E}.
\label{eq:ML}
\end{equation}

Note that direct substitution of the Landauer \emph{lower bound} into the \emph{denominator} of \eqref{eq:ML} does not yield a valid lower bound on $\Delta\tau$: a system consuming more energy than the Landauer minimum may tick \emph{faster}. To obtain a definite relation between time and entropy, we model a \emph{thermodynamically optimal} Witness---one that processes information at the Landauer limit, so that the energy consumed per relational step is exactly the erasure cost:
\begin{equation}
\Delta E = k_B T \ln 2 \;\cdot\; \Delta S(\rho_W).
\label{eq:optimal}
\end{equation}

Substituting \eqref{eq:optimal} into \eqref{eq:ML} yields the relational time step for an optimal Witness:

\begin{axiom}[Relational Time Step for an Optimal Witness]
\begin{equation}
\boxed{\;\Delta \tau_{\mathrm{opt}} \geq \frac{\pi \hbar}{2\, k_B T \ln 2 \;\cdot\; \Delta S(\rho_W)}\;}
\label{eq:taumin}
\end{equation}
\end{axiom}

\begin{remark}[Physical interpretation]
\label{rem:ceiling}
$\Delta\tau_{\mathrm{opt}}$ characterizes the \emph{slowest} permissible clock for a given entropy increment, not the fastest. A thermodynamically optimal Witness operates at the Landauer limit---the quasistatic regime of maximum efficiency and minimum energy expenditure. Any real (non-optimal) Witness dissipates $\Delta E > k_B T \ln 2 \cdot \Delta S$, and the Margolus--Levitin bound permits it to tick \emph{faster} than $\Delta\tau_{\mathrm{opt}}$. The bound is therefore a \emph{ceiling} on the time step (an upper bound on how slowly a witness can tick per bit of entropy), not a floor. $\Delta\tau_{\mathrm{opt}}$ is the maximum temporal cost per unit of relational entropy generation.
\end{remark}

\begin{remark}[The energy identification]
\label{rem:energy}
This derivation identifies two physically distinct quantities: the \emph{dynamical energy} $\langle \hat{H} \rangle - E_0$ driving unitary evolution in the Margolus--Levitin bound, and the \emph{dissipated heat} $k_B T \ln 2 \cdot \Delta S$ of irreversible information erasure in the Landauer bound. The identification $\Delta E_{\mathrm{ML}} = \Delta E_{\mathrm{Landauer}}$ models a Witness whose entire available dynamical energy is consumed by information erasure---a maximally dissipative, thermodynamically reversible limit. This is a strong modeling assumption, not a derivation from first principles. The resulting bound is physically suggestive and dimensionally correct, but the strict equivalence between distinguishable state evolution and information erasure remains heuristic in the general case.
\end{remark}

\begin{remark}
Equation~\eqref{eq:taumin} entails that $\Delta\tau_{\mathrm{opt}}$ is inversely proportional to both temperature and entropy generation. Hotter systems with greater entropy production have smaller $\Delta\tau_{\mathrm{opt}}$---meaning the ceiling on their time step is lower, permitting faster relational dynamics. Conversely, as $\Delta S \to 0$ (a nearly isolated Witness), $\Delta\tau_{\mathrm{opt}} \to \infty$: the maximum time step diverges, and relational time freezes. This is consistent with the global stasis of the Wheeler--DeWitt equation (Axiom~I)---a system with no entropy generation has no relational time.
\end{remark}

\begin{remark}
This result is a quantum-information-theoretic analogue of the thermal time hypothesis of Connes and Rovelli \cite{Connes1994}, which proposes that the flow of time in generally covariant quantum theories emerges from the modular automorphism of the thermal state. Both frameworks anchor temporal flow to thermodynamic structure; the present derivation does so via information-erasure bounds rather than modular theory.
\end{remark}

\noindent Here $\Delta S(\rho_W)$ denotes the incremental change in entanglement entropy per relational step---the number of bits of information about $R$ that become inaccessible to $W$ in the transition. The temperature $T$ appearing in \eqref{eq:taumin} requires careful justification: because the global state is pure ($\hat{H}\ket{\Psi} = 0$), the reduced state $\rho_W$ does not automatically carry a physical temperature. Two routes are available:

\medskip
\noindent\textbf{The modular route.} Define $T$ via the modular Hamiltonian $K_W = -\ln \rho_W$, following the algebraic framework of Connes and Rovelli \cite{Connes1994}. This definition is intrinsic to $\rho_W$ and requires no assumption about thermalization. However, the KMS temperature associated with the modular automorphism is $\beta = 1$ by construction in modular time; it acquires nontrivial physical content only when compared to an independent time flow. Since time is itself defined by Axiom~IV using $T$, there is a circularity risk: the modular temperature requires a time flow for physical interpretation, but time requires $T$. This circularity is the same one that lives within the Connes--Rovelli thermal time hypothesis itself. Additionally, the modular Hamiltonian requires $\rho_W$ to be faithful (no zero eigenvalues), which may not hold for all physically relevant states.

\medskip
\noindent\textbf{The ETH route.} Invoke the Eigenstate Thermalization Hypothesis (ETH) \cite{Srednicki1994}, under which $R$ acts as a thermalizing bath for $W$ and the reduced state $\rho_W$ is approximately thermal at a well-defined physical temperature. This avoids the modular circularity by grounding $T$ in the spectral properties of $\hat{H}$ restricted to the relevant energy shell, but introduces an additional physical assumption about the thermalization properties of the global Hamiltonian.

\medskip
\noindent Both routes yield a consistent framework. We regard the choice between them as open and note that the framework's structural conclusions---the coupling between entropy generation and relational time---hold under either interpretation of $T$.

% ====================================================================
\section{Axiom V: Locating the Hard Problem}
\label{sec:axiom5}

We possess a physical system generating epistemic entropy ($\Delta S > 0$) and a relational parameter tracking it ($\tau$). However, this does not explain \emph{why} $\tau$ is phenomenologically experienced as a flowing ``now.''

The contribution of this framework is not to solve the hard problem of consciousness, but to \emph{locate it with mathematical precision}---to define the structural constraints that any future solution must satisfy.

\begin{axiom}[The Functor Constraint]
Let $\mathbf{Phys}$ be the category whose objects are entangled quantum states $\rho_W$ (with $E_{\mathcal{N}} > 0$) and whose morphisms are local operations and classical communication (LOCC). Let $\mathbf{Phen}$ be the (as yet formally undefined) category of phenomenal experience. The hard problem of consciousness is the joint problem of:
\begin{enumerate}[label=(\roman*)]
    \item defining the category $\mathbf{Phen}$---its objects (experiential states) and morphisms (transitions or relations between experiences); and
    \item constructing a structure-preserving functor $F$ mapping the physical generation of entropy into the subjective generation of experience:
\end{enumerate}
\begin{equation}
F : \mathbf{Phys}(\Delta S > 0) \longrightarrow \mathbf{Phen}(\text{``Now''}).
\label{eq:functor}
\end{equation}
The framework does not presume that only $F$ remains to be defined; the formalization of $\mathbf{Phen}$ is equally (and perhaps more fundamentally) open. What the framework provides is the constraint that $F$ must be a functor---it must preserve the compositional structure of LOCC morphisms in $\mathbf{Phys}$---and that it must factor through the entanglement topology of $W$ (Section~\ref{sec:topo}).
\end{axiom}

While Integrated Information Theory (IIT) \cite{Tononi2016} attempts to quantify the physical correlates of consciousness via the integrated information $\Phi$ of a system's causal architecture, the Relational Witnessing Framework proposes a different structural approach: that $F$ factors through a topological invariant of the entanglement structure internal to $W$.

\begin{remark}[Necessary but not sufficient]
Non-zero entanglement entropy ($E_{\mathcal{N}}(W{:}R) > 0$) is a formally \emph{necessary} condition for witnessing within this framework, but it is not claimed to be \emph{sufficient} for phenomenal experience. A single qubit entangled with a photon satisfies Axiom~III but is not thereby conscious. The emergence of phenomenal experience likely requires sufficient topological complexity within the entanglement complex $\mathcal{K}_W$---for instance, nontrivial persistent homology in dimension $k \geq 1$ (Section~\ref{sec:topo}). The framework constrains the form that any sufficiency condition must take, but does not yet classify the threshold of topological complexity required.
\end{remark}

% ====================================================================
\section{The Topological Hypothesis: Formalizing $\mathcal{G}_W$}
\label{sec:topo}

We now define the entanglement graph $\mathcal{G}_W$ and its topological invariants. This section is conjectural and constitutes a research program, not a derivation.

\medskip
\noindent\textbf{Relation to prior work.} The construction of simplicial complexes from multipartite quantum states, with pairwise entanglement or correlation measures as filtration parameters, is an established direction in topological data analysis. Di~Pierro, Mancini, Memarzadeh, and Mengoni \cite{DiPierro2018} first proposed the use of bipartite entanglement as a distance for Vietoris--Rips filtrations on multipartite qubit states; Mengoni, Memarzadeh, Di~Pierro, and Mancini \cite{Mengoni2019} developed the approach further. Olsthoorn \cite{Olsthoorn2023} constructed filtered simplicial complexes from inverse quantum mutual information and applied persistent homology to condensed-matter spin systems. Hamilton and Leditzky \cite{HamiltonLeditzky2024} refined the construction using \v{C}ech complexes weighted by Tsallis-deformed total correlation, defining an Integrated Euler Characteristic as a topological summary of multipartite entanglement; operational interpretations of this summary via distillable-entanglement bounds have subsequently been developed \cite{Spellman2025}. The construction given below specializes this general program to logarithmic negativity as the pairwise weight (motivated by Axiom~III, where mixedness of the reduced state is the generic case), adopts the flag complex for its face-closure property and computational tractability, and embeds the resulting signature within the broader categorical framework of Axiom~V---specifically, as the invariant through which the functor $F$ is hypothesized to factor (Section~\ref{sec:axiom5}). The application domain (relational time and phenomenal experience) is also distinct: prior work has studied abstract multipartite states and condensed-matter spin systems, not the use of entanglement topology as a constraint on theories of subjective experience.

% ---- Step 1: Internal Factorization ----
\subsection{Internal Factorization of $W$}

For the entanglement structure internal to $W$ to be well-defined, $\mathcal{H}_W$ must factor into subsystems:
\begin{equation}
\mathcal{H}_W = \mathcal{H}_1 \otimes \mathcal{H}_2 \otimes \cdots \otimes \mathcal{H}_n.
\label{eq:internal}
\end{equation}

The same factorization-dependence issue identified by Zanardi et al.\ \cite{Zanardi2004} applies here. We resolve it by recursive application of the algebraic criterion from Section~\ref{sec:axiom2}: the preferred internal factorization is the one with respect to which the Witness's internal observables act locally. In the emergent-time regime, this coincides with the sub-factorization stabilized by internal decoherence, as in Zurek's einselection program \cite{Zurek2009}.

% ---- Step 2: Weighted Simplicial Complex ----
\subsection{The Entanglement Complex $\mathcal{K}_W$}

\begin{definition}[Entanglement Complex]
\label{def:KW}
Given the internal factorization \eqref{eq:internal} and the reduced state $\rho_W$, we define the \emph{entanglement complex} $\mathcal{K}_W$ as a \emph{flag complex} (clique complex) constructed from the weighted 1-skeleton of pairwise entanglement:
\begin{itemize}
    \item \textbf{0-simplices} (vertices): individual subsystems $\{1, 2, \ldots, n\}$.
    \item \textbf{1-simplices} (edges): pairs $(i,j)$ with weight given by the pairwise logarithmic negativity
    \begin{equation}
    w_{ij} = E_{\mathcal{N}}(i{:}j) = \log_2 \lVert \rho_{ij}^{T_i} \rVert_1,
    \label{eq:negativity}
    \end{equation}
    where $\rho_{ij} = \operatorname{Tr}_{\overline{ij}}(\rho_W)$ is the reduced state of the pair and $T_i$ denotes partial transpose with respect to subsystem $i$.
    \item \textbf{$k$-simplices} ($k \geq 2$): a subset $\{i_1, \ldots, i_{k+1}\}$ is included as a $k$-simplex if and only if \emph{all} $\binom{k+1}{2}$ pairwise edges $(i_a, i_b)$ are present (i.e., have $E_{\mathcal{N}}(i_a{:}i_b) > 0$).
\end{itemize}
\end{definition}

The choice of logarithmic negativity is consistent with Axiom~III: it detects genuine quantum entanglement and is zero for separable states. The flag complex construction is adopted for two reasons. First, it guarantees the face-closure property required of a valid simplicial complex: every face of an included simplex is itself included. An alternative construction weighting higher simplices by genuine multipartite entanglement measures (e.g., the $n$-tangle \cite{Coffman2000} or genuine multipartite negativity \cite{Jungnitsch2011}) can violate this property---a GHZ state of three qubits, for instance, possesses maximum tripartite entanglement but zero pairwise entanglement after partial trace, yielding a 2-simplex without its 1-simplex faces. Second, the flag complex is entirely determined by the 1-skeleton and is computationally tractable, since computing genuine multipartite entanglement measures for $k \geq 2$ is in general NP-hard for large systems.

\begin{remark}
The flag complex may fail to capture genuinely multipartite entanglement that exists independently of pairwise correlations; this limitation is intrinsic to any construction derived from pair marginals and is well-known in the broader persistent-homology-of-entanglement literature \cite{DiPierro2018, HamiltonLeditzky2024}. We regard a richer simplicial construction that accommodates such states---while preserving face-closure---as an important open problem in the formalization of $\mathcal{K}_W$.
\end{remark}

% ---- Step 3: Persistent Homology ----
\subsection{Topological Invariants via Persistent Homology}

Rather than choosing an arbitrary threshold to binarize the weighted complex, we apply \emph{persistent homology} \cite{Edelsbrunner2002, Zomorodian2005}, which tracks how topological features---connected components ($H_0$), loops ($H_1$), voids ($H_2$), and higher-dimensional cavities ($H_k$)---appear and disappear as a filtration parameter $\varepsilon$ sweeps from the maximum weight to zero.

At threshold $\varepsilon$, include all simplices $\sigma \in \mathcal{K}_W$ with weight $w(\sigma) \geq \varepsilon$. As $\varepsilon$ decreases, more simplices enter the complex, and the homology changes. Each topological feature is born at some $\varepsilon_b$ and dies at some $\varepsilon_d < \varepsilon_b$, yielding an interval $[\varepsilon_d, \varepsilon_b)$.

\begin{definition}[Topological Signature]
The \emph{topological signature} of the Witness is the collection of persistence diagrams across all homological dimensions:
\begin{equation}
\chi(\mathcal{K}_W) = \bigl\{ \mathrm{Dgm}_k(\mathcal{K}_W) \bigr\}_{k=0}^{\dim \mathcal{K}_W},
\label{eq:persistence}
\end{equation}
where each $\mathrm{Dgm}_k$ is a multiset of intervals $[\varepsilon_d, \varepsilon_b)$ recording the birth and death of $k$-dimensional features.
\end{definition}

Persistent homology has the desired properties for our framework:
\begin{enumerate}[label=(\roman*)]
    \item \textbf{Stability:} small perturbations in the entanglement weights produce small changes in the persistence diagrams (with respect to the bottleneck or Wasserstein distance \cite{CohenSteiner2007}). This ensures robustness under minor fluctuations.
    \item \textbf{Sensitivity to structure:} qualitative changes in the entanglement topology (e.g., severing a cluster of connections, as under anesthesia) produce measurable changes in the persistence diagrams.
    \item \textbf{Computability:} persistent homology is algorithmically tractable for finite complexes.
\end{enumerate}

% ---- Step 4: The Fiber Bundle Structure ----
\subsection{Class vs.\ Content: The Fiber Bundle}

The persistent homology framework naturally formalizes the distinction between the \emph{class} of experience and its \emph{content}.

Let $\mathcal{B}$ be the space of persistence diagrams, equipped with the Wasserstein metric $d_W$. For each point $b \in \mathcal{B}$, define the \emph{fiber}:
\begin{equation}
\mathcal{F}_b = \bigl\{ \rho_W \;\big|\; \chi(\mathcal{K}_W[\rho_W]) = b \bigr\},
\label{eq:fiber}
\end{equation}
the set of all Witness states whose entanglement complex has topological signature $b$.

\begin{proposition}[Two-Level Structure]
The functor $F$ decomposes as:
\begin{equation}
F(\rho_W) = F'\!\bigl(\chi(\mathcal{K}_W),\; \rho_W\bigr),
\label{eq:twolevel}
\end{equation}
where:
\begin{itemize}
    \item $\chi(\mathcal{K}_W) \in \mathcal{B}$ determines the \emph{structural class} of experience (e.g., wakefulness vs.\ dreamless sleep vs.\ anesthesia);
    \item the specific state $\rho_W$ within the fiber $\mathcal{F}_{\chi(\mathcal{K}_W)}$ determines the \emph{phenomenal content} (e.g., seeing red vs.\ hearing a tone).
\end{itemize}
\end{proposition}

This structure is that of a surjection $\pi: \mathcal{E} \to \mathcal{B}$ with level sets, where $\mathcal{E}$ is the total space of Witness states, $\mathcal{B}$ is the base space of topological signatures, and each $\mathcal{F}_b = \pi^{-1}(b)$ is the set of states sharing that signature. The map $\rho_W \mapsto \chi(\mathcal{K}_W)$ is continuous by the stability theorem for persistent homology \cite{CohenSteiner2007}, but whether the level sets $\mathcal{F}_b$ are locally homeomorphic to each other---the condition required for a genuine fiber bundle---is an open question. If local triviality holds, $\pi$ acquires the full structure of a fiber bundle; we conjecture but do not prove this.

\medskip
\noindent The Topological Hypothesis, as developed through Sections 6.1--6.4, provides the following: a well-defined simplicial complex $\mathcal{K}_W$ encoding the entanglement structure of the Witness, a stable topological invariant $\chi(\mathcal{K}_W)$ via persistent homology, and a two-level decomposition of the functor $F$ into structural class and phenomenal content. Whether the space of topological signatures $\mathcal{B}$ itself has interesting higher homotopy structure---and in particular whether such structure might exhibit periodicity---is an open problem explored in Appendix~\ref{app:bott}.

% ====================================================================
\section{Structural Consequences and Empirical Constraints}
\label{sec:empirical}

We outline three structural consequences of the framework. The first two are interpretive---they follow from the formalism but do not predict deviations from standard quantum mechanics. The third is in principle empirically distinguishable.

\begin{enumerate}
    \item \textbf{Relational time freezing.} For a perfectly isolated quantum system ($\Delta S(\rho_W) \to 0$), Equation~\eqref{eq:taumin} yields $\Delta\tau_{\mathrm{opt}} \to \infty$: relational time ceases. This is not a ``temporal anomaly'' measurable against an external clock; it is the framework's statement that a system generating no entanglement entropy has no relational time parameter at all. This is consistent with the global stasis of the Wheeler--DeWitt equation (Axiom~I) and with the Page--Wootters mechanism, in which time is strictly relational. We note that this consequence is \emph{interpretive}: it does not predict any observable deviation from the Schr\"{o}dinger equation's predictions for isolated systems. The framework is empirically equivalent to standard quantum mechanics in this regime; the difference is ontological, not operational.
    
    \item \textbf{The entanglement prerequisite for experience.} A system completely isolated into a pure state ($S = 0$, $E_{\mathcal{N}} = 0$) cannot, by definition, possess a temporal ``now'' or act as a formal Witness. Genuine quantum entanglement is a strict physical prerequisite for phenomenal experience within this framework.
    
    \item \textbf{Topological transitions under anesthesia.} If the class of experience is determined by the persistence diagrams $\chi(\mathcal{K}_W)$, then the transition from wakefulness to anesthesia-induced unconsciousness should correspond to a measurable topological transition in the entanglement structure---specifically, a collapse in the persistent homology (loss of long-lived topological features in $H_1$ and higher). This is in principle testable via quantum-information-theoretic analysis of neural correlates, extending the existing work connecting neural complexity measures to consciousness \cite{Casali2013}.
\end{enumerate}

% ====================================================================
\section{The Synthesized Formalism}
\label{sec:synthesis}

The Relational Witnessing Framework proposes that witnessing, entropy, and time are three descriptions of a single relational structure. We now make this claim precise. The three concepts are \emph{co-generated by a single mathematical operation}: the partial trace $\operatorname{Tr}_R$ applied to the bipartition $\mathcal{H}_W \otimes \mathcal{H}_R$. This operation simultaneously:
\begin{itemize}
    \item defines the epistemic boundary of the Witness (what the subsystem can access);
    \item produces the mixed state $\rho_W$ whose von Neumann entropy $S(\rho_W) > 0$ is the entanglement entropy; and
    \item via conditionalization on clock states $\ket{w(\tau)}$, yields the relational time parameter $\tau$.
\end{itemize}

\noindent The identity claim is that these are not three independently motivated concepts that happen to coincide, but three projections of the structure induced by a single operation (the epistemic cut). A stronger formulation would require proving that witnessing, entropy, and time are \emph{isomorphic} as objects in some encompassing category---that there exists a formal sense in which any one determines the other two. We regard this stronger claim as a conjecture motivated by the co-generation structure, not a proven theorem. The synthesized formalism is:

\begin{align}
\hat{H}\ket{\Psi} &= 0 &&\text{(The Timeless Whole)} \label{eq:syn1}\\[6pt]
\mathcal{H}_{\mathrm{total}} &= \mathcal{H}_W \otimes \mathcal{H}_R &&\text{(The Algebraic Bipartition)} \label{eq:syn2}\\[6pt]
\rho_W &= \operatorname{Tr}_R\!\bigl(\ket{\Psi}\!\bra{\Psi}\bigr),\quad E_{\mathcal{N}}(W{:}R) > 0 &&\text{(The Epistemic Cut)} \label{eq:syn3}\\[6pt]
\Delta\tau_{\mathrm{opt}} &\geq \frac{\pi\hbar}{2\,k_BT\ln 2 \;\cdot\; \Delta S(\rho_W)} &&\text{(Relational Time)} \label{eq:syn4}\\[6pt]
F(\rho_W) &= F'\!\bigl(\chi(\mathcal{K}_W),\; \rho_W\bigr) &&\text{(The Topological Constraint)} \label{eq:syn5}
\end{align}

% ====================================================================
\section{Conclusion}

This framework makes a specific structural claim: that time, entropy, and witnessing are not three separate phenomena requiring three separate explanations, but three projections of a single relational structure induced by the partial trace across a bipartition of a timeless whole. The claim of co-generation is demonstrated; the stronger claim of formal isomorphism remains an open conjecture. The deductive core (Axioms~I--IV) rests on established physics: the Wheeler--DeWitt equation, the Page--Wootters mechanism, Landauer's principle, and the Margolus--Levitin speed limit, composed under an explicit thermodynamic optimality assumption. The conjectural extension (Axiom~V and the Topological Hypothesis) proposes that the hard problem of consciousness can be located---if not yet solved---as the joint problem of defining both the category of phenomenal experience and a structure-preserving functor from the category of entangled quantum states into it, and that this functor factors through the persistent homology of the entanglement complex.

We acknowledge the most dangerous foundational objection directly: the framework's conclusions depend on the choice of tensor product structure (the bipartition), and different factorizations yield different witnesses, entropies, and time parameters. If the preferred factorization is not uniquely determined, the framework's outputs may be artifacts of a decomposition choice rather than features of reality. We have offered an algebraic criterion (observable-induced factorization) supplemented by a dynamical consistency check (Quantum Darwinism), but whether this fully resolves the ambiguity is an open question. The framework stands or falls on the defensibility of its factorization criterion.

The framework connects naturally to several active research programs: the quantum reference frame formalism of H\"{o}hn, Smith, and Lock \cite{Hohn2021}, which provides a rigorous foundation for the Page--Wootters mechanism; the thermal time hypothesis of Connes and Rovelli \cite{Connes1994}, which similarly anchors temporal flow to thermodynamic structure; the ER\,=\,EPR program of Maldacena and Susskind \cite{Maldacena2013} and Van~Raamsdonk's entanglement--geometry correspondence \cite{VanRaamsdonk2010}, which suggest that entanglement entropy generates spatial geometry as well as time; and topological entanglement entropy in the sense of Kitaev and Preskill \cite{Kitaev2006}, which provides precedent for extracting topological invariants from entanglement structure.

Whether the functor $F$ can be explicitly constructed, and whether the space of topological signatures exhibits further structure (see Appendix~\ref{app:bott}), are open mathematical questions. What the Relational Witnessing Framework provides is the precise scaffold---the categorical, topological, and information-theoretic constraints---that any such construction must satisfy.

% ====================================================================
\appendix
\section{Open Problem: Periodicity in the Entanglement Topology}
\label{app:bott}

This appendix poses a speculative but precise open problem motivated by the topological framework of Section~\ref{sec:topo}. It is structurally independent of the main results: the framework's deductive core (Axioms~I--IV) and the Topological Hypothesis (Section~6, through the fiber structure of Section~6.4) do not depend on the conjecture stated here.

\medskip

The mathematician Raoul Bott identified a deep periodicity in the homotopy groups of the classical groups: $\pi_{k+8}(O) \cong \pi_k(O)$ and $\pi_{k+2}(U) \cong \pi_k(U)$ \cite{Bott1959}. This periodicity governs the classification of fiber bundles and the structure of $K$-theory. We ask whether an analogous periodicity appears in the topology of entanglement.

\begin{conjecture}[Bott Periodicity in Entanglement Topology]
\label{conj:bott}
Let $\mathcal{B}_n$ denote the space of topological signatures $\chi(\mathcal{K}_W)$ for Witnesses with $n$ internal subsystems. As $n$ increases, the homotopy groups $\pi_k(\mathcal{B}_n)$ exhibit periodicity:
\begin{equation}
\pi_{k+p}(\mathcal{B}_n) \cong \pi_k(\mathcal{B}_n) \quad \text{for some fixed period } p,
\label{eq:bottconj}
\end{equation}
reflecting a fundamental periodicity in the structural classes of entanglement topology available to witnessing systems of increasing complexity.
\end{conjecture}

If true, this would imply that the ``menu'' of structurally distinct experience-classes is not unbounded but recurs with a characteristic period---that beyond a certain complexity, new subsystems do not generate fundamentally new topological classes, but repeat existing ones at higher resolution.

The natural route to establishing Conjecture~\ref{conj:bott} would be to show that the entanglement complexes $\mathcal{K}_W$, equipped with the local Hilbert spaces $\mathcal{H}_i$ as fibers over their vertices, define complex vector bundles whose classifying space is $BU(n)$---the classifying space of the unitary group, where Bott periodicity natively lives ($\pi_{k+2}(BU) \cong \pi_k(BU)$). Analogous applications of $K$-theoretic periodicity to physical quantum systems have been developed, most notably in Kitaev's periodic-table classification of topological insulators and superconductors \cite{Kitaev2009}, which grounds the plausibility---though not the validity---of seeking such periodicity in the space of entanglement signatures. If a map from entanglement complexes to $BU$ can be constructed, the periodicity of $\pi_k(\mathcal{B}_n)$ would follow from the periodicity of $\pi_k(BU)$ via the induced homomorphism on homotopy groups.

Two significant obstacles must be acknowledged:

\begin{enumerate}
    \item \textbf{Transition functions.} A collection of vector spaces over vertices does not constitute a vector bundle without \emph{transition functions} along the edges that glue the fibers together. To rigorously map $\mathcal{K}_W$ to $BU(n)$, future work must define these transition functions---potentially utilizing the local modular flow or the entanglement Hamiltonian $K_{ij} = -\ln \rho_{ij}$ to provide the requisite holonomy between $\mathcal{H}_i$ and $\mathcal{H}_j$ along each 1-simplex. Without such transition functions, the bundle is trivial and the conjecture is void.
    
    \item \textbf{Contractibility.} The space of persistence diagrams equipped with the Wasserstein or bottleneck metric is known to be contractible for certain natural filtrations \cite{Mileyko2011}. If $\mathcal{B}_n$ is contractible, then $\pi_k(\mathcal{B}_n) = 0$ for all $k \geq 1$, and the periodicity conjecture is trivially satisfied ($0 \cong 0$) but content-free. The conjecture has nontrivial content only if the relevant subspace of $\mathcal{B}_n$---the persistence diagrams \emph{actually realizable} by entanglement complexes of $n$-partite quantum states---is non-contractible. Determining the homotopy type of this physically realizable subspace is an open problem at the intersection of topological data analysis, $K$-theory, and quantum information theory.
\end{enumerate}

We state this conjecture because it is precise, falsifiable, and---if true---would have profound consequences for the structure of the functor $F$. But we separate it from the main body of the paper to make clear that the Relational Witnessing Framework does not depend on it.

% ====================================================================
\begin{thebibliography}{99}

\bibitem{DeWitt1967}
B.~S.~DeWitt, ``Quantum Theory of Gravity. I. The Canonical Theory,'' \textit{Phys.\ Rev.}\ \textbf{160}, 1113 (1967).

\bibitem{Zanardi2004}
P.~Zanardi, D.~A.~Lidar, and S.~Lloyd, ``Quantum tensor product structures are observable induced,'' \textit{Phys.\ Rev.\ Lett.}\ \textbf{92}, 060402 (2004).

\bibitem{Zurek2009}
W.~H.~Zurek, ``Quantum Darwinism,'' \textit{Nature Phys.}\ \textbf{5}, 181 (2009).

\bibitem{Page1983}
D.~N.~Page and W.~K.~Wootters, ``Evolution without evolution: Dynamics described by stationary observables,'' \textit{Phys.\ Rev.\ D}\ \textbf{27}, 2885 (1983).

\bibitem{Landauer1961}
R.~Landauer, ``Irreversibility and Heat Generation in the Computing Process,'' \textit{IBM J.\ Res.\ Dev.}\ \textbf{5}, 183 (1961).

\bibitem{Margolus1998}
N.~Margolus and L.~B.~Levitin, ``The maximum speed of dynamical evolution,'' \textit{Physica D}\ \textbf{120}, 188 (1998).

\bibitem{Connes1994}
A.~Connes and C.~Rovelli, ``Von Neumann algebra automorphisms and time-thermodynamics relation in generally covariant quantum theories,'' \textit{Class.\ Quantum Grav.}\ \textbf{11}, 2899 (1994).

\bibitem{Tononi2016}
G.~Tononi, M.~Boly, M.~Massimini, and C.~Koch, ``Integrated information theory: from consciousness to its physical substrate,'' \textit{Nature Rev.\ Neurosci.}\ \textbf{17}, 450 (2016).

\bibitem{Coffman2000}
V.~Coffman, J.~Kundu, and W.~K.~Wootters, ``Distributed entanglement,'' \textit{Phys.\ Rev.\ A}\ \textbf{61}, 052306 (2000).

\bibitem{Jungnitsch2011}
B.~Jungnitsch, T.~Moroder, and O.~G\"{u}hne, ``Taming multiparticle entanglement,'' \textit{Phys.\ Rev.\ Lett.}\ \textbf{106}, 190502 (2011).

\bibitem{Edelsbrunner2002}
H.~Edelsbrunner, D.~Letscher, and A.~Zomorodian, ``Topological persistence and simplification,'' \textit{Discrete Comput.\ Geom.}\ \textbf{28}, 511 (2002).

\bibitem{Zomorodian2005}
A.~Zomorodian and G.~Carlsson, ``Computing persistent homology,'' \textit{Discrete Comput.\ Geom.}\ \textbf{33}, 249 (2005).

\bibitem{CohenSteiner2007}
D.~Cohen-Steiner, H.~Edelsbrunner, and J.~Harer, ``Stability of persistence diagrams,'' \textit{Discrete Comput.\ Geom.}\ \textbf{37}, 103 (2007).

\bibitem{DiPierro2018}
A.~Di~Pierro, S.~Mancini, L.~Memarzadeh, and R.~Mengoni, ``Homological analysis of multi-qubit entanglement,'' \textit{Europhys.\ Lett.}\ \textbf{123}, 30006 (2018), arXiv:1802.04572.

\bibitem{Mengoni2019}
R.~Mengoni, L.~Memarzadeh, A.~Di~Pierro, and S.~Mancini, ``Persistent homology analysis of multiqubit entanglement,'' arXiv:1907.06914 (2019).

\bibitem{Olsthoorn2023}
B.~Olsthoorn, ``Persistent homology of quantum entanglement,'' \textit{Phys.\ Rev.\ B}\ \textbf{107}, 115174 (2023), arXiv:2110.10214.

\bibitem{HamiltonLeditzky2024}
G.~A.~Hamilton and F.~Leditzky, ``Probing multipartite entanglement through persistent homology,'' \textit{Commun.\ Math.\ Phys.}\ \textbf{405}, 125 (2024), arXiv:2307.07492.

\bibitem{Spellman2025}
C.~Spellman \textit{et al.}, ``Interpreting Multipartite Entanglement through Topological Summaries,'' arXiv:2505.00642 (2025).

\bibitem{Bott1959}
R.~Bott, ``The stable homotopy of the classical groups,'' \textit{Ann.\ of Math.}\ \textbf{70}, 313 (1959).

\bibitem{Woods2019}
M.~P.~Woods, R.~Silva, and J.~Oppenheim, ``Autonomous quantum machines and finite-sized clocks,'' \textit{Ann.\ Henri Poincar\'{e}}\ \textbf{20}, 125 (2019).

\bibitem{Casali2013}
A.~G.~Casali \textit{et al.}, ``A theoretically based index of consciousness independent of sensory processing and behavior,'' \textit{Sci.\ Transl.\ Med.}\ \textbf{5}, 198ra105 (2013).

\bibitem{Hohn2021}
P.~A.~H\"{o}hn, A.~R.~H.~Smith, and M.~P.~E.~Lock, ``Trinity of relational quantum dynamics,'' \textit{Phys.\ Rev.\ D}\ \textbf{104}, 066001 (2021).

\bibitem{Maldacena2013}
J.~Maldacena and L.~Susskind, ``Cool horizons for entangled black holes,'' \textit{Fortschr.\ Phys.}\ \textbf{61}, 781 (2013).

\bibitem{VanRaamsdonk2010}
M.~Van~Raamsdonk, ``Building up spacetime with quantum entanglement,'' \textit{Gen.\ Rel.\ Grav.}\ \textbf{42}, 2323 (2010).

\bibitem{Kitaev2006}
A.~Kitaev and J.~Preskill, ``Topological entanglement entropy,'' \textit{Phys.\ Rev.\ Lett.}\ \textbf{96}, 110404 (2006).

\bibitem{Kitaev2009}
A.~Kitaev, ``Periodic table for topological insulators and superconductors,'' \textit{AIP Conf.\ Proc.}\ \textbf{1134}, 22 (2009), arXiv:0901.2686.

\bibitem{Srednicki1994}
M.~Srednicki, ``Chaos and quantum thermalization,'' \textit{Phys.\ Rev.\ E}\ \textbf{50}, 888 (1994).

\bibitem{Mileyko2011}
Y.~Mileyko, S.~Mukherjee, and J.~Harer, ``Probability measures on the space of persistence diagrams,'' \textit{Inverse Problems}\ \textbf{27}, 124007 (2011).

\end{thebibliography}

\end{document}
